\section{Form HTML dan Pemrosesan Form dengan PHP}

Form HTML adalah mekanisme utama bagi pengguna untuk mengirimkan data ke server web. Setiap form memiliki atribut \code{action} (URL tujuan pengiriman) dan \code{method} (metode HTTP: GET atau POST). Elemen-elemen form seperti \code{<input>}, \code{<textarea>}, dan \code{<select>} memiliki atribut \code{name} yang menjadi kunci di array \code{\$\_GET} atau \code{\$\_POST} saat data diterima oleh PHP. Pemahaman alur pengiriman dan pemrosesan form ini menjadi dasar dari hampir semua aplikasi web interaktif \cite{phpManual}.

Alur pemrosesan form dimulai ketika pengguna mengisi form dan menekan tombol \textit{submit}. Browser mengemas data form dan mengirimkannya ke server sesuai metode yang ditentukan. PHP menerima data tersebut melalui superglobal \code{\$\_GET} atau \code{\$\_POST}, memvalidasinya, dan menghasilkan respons. Validasi server-side (\textit{server-side validation}) \textbf{wajib} dilakukan meskipun sudah ada validasi di sisi klien (HTML5 atau JavaScript), karena validasi klien dapat dengan mudah dilewati oleh pengguna yang mahir \cite{phpRightWay}.

\begin{figure}[ht]
\centering
\begin{tikzpicture}[
  node distance=0.8cm,
  box/.style={rectangle, draw, fill=blue!15, rounded corners, text width=3cm, minimum height=0.9cm, text centered, font=\small},
  decision/.style={diamond, draw, fill=orange!20, text width=2cm, text badly centered, inner sep=2pt, aspect=2, font=\small},
  line/.style={draw, -Stealth, thick},
  startstop/.style={rectangle, draw, fill=green!15, rounded corners, minimum width=2cm, minimum height=0.7cm, text centered, font=\small}
]
  \node[startstop] (start) {User mengisi form};
  \node[box, below=of start] (submit) {Klik Submit\\(POST/GET)};
  \node[box, below=of submit] (receive) {PHP menerima\\\code{\$\_POST}/\code{\$\_GET}};
  \node[decision, below=of receive] (valid) {Data valid?};
  \node[box, right=2.5cm of valid] (error) {Tampilkan pesan error};
  \node[box, below=of valid] (process) {Proses data\\(simpan ke DB, email)};
  \node[startstop, below=of process] (success) {Tampilkan sukses};

  \path[line] (start) -- (submit);
  \path[line] (submit) -- (receive);
  \path[line] (receive) -- (valid);
  \path[line] (valid) -- node[above, font=\scriptsize]{Tidak} (error);
  \path[line] (valid) -- node[right, font=\scriptsize]{Ya} (process);
  \path[line] (error) |- (start);
  \path[line] (process) -- (success);
\end{tikzpicture}
\caption{Alur pemrosesan form dengan PHP}
\end{figure}

Berikut contoh form HTML sederhana yang mengirimkan data dengan metode POST dan file PHP yang memprosesnya. Perhatikan bahwa form dan pemrosesan bisa berada di file yang sama (self-processing form) atau di file terpisah.

\begin{webcode}[language=HTML, caption={Form HTML sederhana}]
<!-- form_kontak.php -->
<!DOCTYPE html>
<html lang="id">
<head><title>Form Kontak</title></head>
<body>
  <h2>Hubungi Kami</h2>
  <form method="post" action="form_kontak.php">
    <label>Nama: <input type="text" name="nama"></label><br><br>
    <label>Email: <input type="email" name="email"></label><br><br>
    <label>Pesan:<br>
      <textarea name="pesan" rows="4" cols="40"></textarea>
    </label><br><br>
    <button type="submit">Kirim</button>
  </form>
</body>
</html>
\end{webcode}

Validasi server-side harus mencakup beberapa pemeriksaan: apakah field wajib sudah diisi, apakah format data benar (misalnya email valid), dan apakah panjang data dalam batas wajar. PHP menyediakan fungsi \code{filter\_var()} untuk validasi dan sanitasi data. Misalnya, \code{filter\_var(\$email, FILTER\_VALIDATE\_EMAIL)} mengembalikan email jika formatnya valid, atau \code{false} jika tidak. Selalu gunakan \code{htmlspecialchars()} saat menampilkan ulang data pengguna untuk mencegah XSS \cite{owaspXss}.

\begin{contoh}
\textbf{Studi Kasus: Form Kontak dengan Validasi Server-Side}

Sebuah website perusahaan membutuhkan form kontak yang menerima nama, email, dan pesan dari pengunjung. Berikut implementasi lengkap dengan validasi server-side menggunakan PHP.

\textbf{Kebutuhan validasi:}
\begin{enumerate}
  \item Semua field wajib diisi (tidak boleh kosong)
  \item Email harus dalam format yang valid
  \item Pesan minimal 10 karakter
  \item Semua output harus disanitasi untuk mencegah XSS
\end{enumerate}
\end{contoh}

\begin{webcode}[language=PHP, caption={Pemrosesan form dengan validasi lengkap}]
<?php
$errors = [];
$sukses = false;

if ($_SERVER["REQUEST_METHOD"] === "POST") {
    // Ambil dan sanitasi input
    $nama  = trim($_POST["nama"] ?? "");
    $email = trim($_POST["email"] ?? "");
    $pesan = trim($_POST["pesan"] ?? "");

    // Validasi: nama tidak boleh kosong
    if ($nama === "") {
        $errors[] = "Nama wajib diisi.";
    }

    // Validasi: email harus valid
    if (!filter_var($email, FILTER_VALIDATE_EMAIL)) {
        $errors[] = "Format email tidak valid.";
    }

    // Validasi: pesan minimal 10 karakter
    if (strlen($pesan) < 10) {
        $errors[] = "Pesan minimal 10 karakter.";
    }

    // Jika tidak ada error, proses data
    if (empty($errors)) {
        // Simpan ke database, kirim email, dll.
        $sukses = true;
    }
}
?>
<!DOCTYPE html>
<html lang="id">
<head><title>Form Kontak</title></head>
<body>
  <h2>Hubungi Kami</h2>

  <?php if ($sukses): ?>
    <p style="color:green;">Pesan Anda berhasil dikirim!</p>
  <?php endif; ?>

  <?php if (!empty($errors)): ?>
    <ul style="color:red;">
      <?php foreach ($errors as $err): ?>
        <li><?= htmlspecialchars($err) ?></li>
      <?php endforeach; ?>
    </ul>
  <?php endif; ?>

  <form method="post">
    <label>Nama:
      <input type="text" name="nama"
             value="<?= htmlspecialchars($nama ?? '') ?>">
    </label><br><br>
    <label>Email:
      <input type="email" name="email"
             value="<?= htmlspecialchars($email ?? '') ?>">
    </label><br><br>
    <label>Pesan:<br>
      <textarea name="pesan" rows="4"
                cols="40"><?= htmlspecialchars($pesan ?? '') ?></textarea>
    </label><br><br>
    <button type="submit">Kirim</button>
  </form>
</body>
</html>
\end{webcode}

\begin{catatan}
\textbf{Teknik \textit{Sticky Form}:} Pada contoh di atas, perhatikan bahwa atribut \code{value} pada input diisi dengan data yang sudah dimasukkan pengguna (menggunakan \code{htmlspecialchars()}). Teknik ini disebut \textit{sticky form} --- jika validasi gagal, pengguna tidak perlu mengisi ulang seluruh form dari awal. Ini meningkatkan pengalaman pengguna secara signifikan \cite{phpRightWay}.
\end{catatan}

